%%%%%%%%%%%%%%%%% Estilo general del Libro %%%%%%%%%%%%%%%%%
\documentclass{wileySev}

%%%%%%%%%%%%%%%%% Algunos paquetes especiales %%%%%%%%%%%%%%%%%
\usepackage[spanish]{babel}
\usepackage[utf8]{inputenc}
\unaccentedoperators
\usepackage[nodayofweek,level]{datetime}
\usepackage{fancyhdr}
\usepackage{longtable}
\usepackage{graphicx}
\usepackage{threeparttablex}
\usepackage{rotating}
\usepackage{fontsize}
  \changefontsize[13.8pt]{12pt}
\usepackage{gensymb}
\usepackage{cancel}
\usepackage{multicol}
\usepackage{tabulary}
\usepackage{nccmath}
\usepackage{chemfig}
\usepackage[hyphens]{url}
\usepackage{float}
\usepackage[
  labelfont=bf,
  singlelinecheck=false,
  indention=2.4cm,
  format=plain,
  parindent=1em,
  justification=centering
]{caption}
\usepackage{subcaption}
\usepackage{tikz}
\usepackage{array,ragged2e, tabularx}
\usepackage{tcolorbox}
\usepackage[export]{adjustbox}
\usepackage{microtype}
\usepackage{wasysym}
\usepackage{colortbl}
\usepackage{hhline}
\usepackage{lscape}
\usetikzlibrary{mindmap}
\usepackage{enumerate} 
\usepackage{enumitem} 
\usepackage{mathtools}
\usepackage{amssymb}
\usepackage{multirow}
\usepackage{etoolbox}
\usepackage{marvosym}
\usepackage[section]{placeins}
\usepackage{calc} 
\usepackage{newfloat}
\usepackage{lipsum}  
\usepackage{xcolor}  
\newcolumntype{P}[1]{>{\raggedRight}p{#1}}

%%%%%%%%%%%%%%%%%% Bibliografía %%%%%%%%%%%%%%%%%%%%%%%%%%%%%

%Estilo de la bibliografía. Puede cambiarlo al estilo IEEE colocando style=ieee
\usepackage[style=apa]{biblatex}

%El archivo bibliografia.bib contiene las entradas de la bibliografia
\addbibresource{bibliografia.bib}

%%%%%%%%%%%%%%%%% Configuraciones especiales %%%%%%%%%%%%%%%%%

\renewcommand\spanishtablename{Tabla}
\tcbuselibrary{skins,breakable}
\usetikzlibrary{shadings,shadows}
\newenvironment{problem}[1]{%
    \tcolorbox[beamer,%
    noparskip,%breakable,
%     colback=LightBlue,colframe=DarkBlue,%
%     colbacklower=DarkBlue!75!LightBlue,%
    title=#1]}%
    {\endtcolorbox}
%%%%%%%%%%%%%%%%%

\addto\captionsspanish{% Reemplaza ''inglés'' con el idioma que usas
  \renewcommand{\contentsname}%
    {ÍNDICE GENERAL}%
}
%%%%%%%%%%%%%%%%%
\definecolor{azul}{rgb}{0.07, 0.04, 0.56}

%%%%%%%%%%%%%%%%%
\makeatletter
\def\hlinewd#1{%
\noalign{\ifnum0=`}\fi\hrule \@height #1 %
\futurelet\reserved@a\@xhline} 
\newcommand{\abs}[1]{\left\vert#1\right\vert}

%%%%%%%%%%%%%%%%%
\apptocmd{\lim}{\limits}{}{}
\DeclarePairedDelimiter\ceil{\lceil}{\rceil}
\DeclarePairedDelimiter\floor{\lfloor}{\rfloor}
\newenvironment{whereparam*}
  {\par\vspace{\abovedisplayskip}\noindent
   \tabularx{\columnwidth}{>{$}l<{$} @{${}={}$} >{\raggedright\arraybackslash}X}}
  {\endtabularx\par\vspace{\belowdisplayskip}}
\addto\captionsspanish{\renewcommand{\chaptername}{Capítulo}}

%%%%%%%%%%%%%%%%% Para texto PostScript
\usepackage{w-bookps}


%%%%%%%%%%%%%%%%%
\def\helveticabold{cmssbx10}

%%%%%%%%%%%%%%%%%
\newcommand{\VT}[1]{\ensuremath{{V_{T#1}}}}

%%%%%%%%%%%%%%%%% Use un cuadro para expandir la macro antes de colocar el encabezado de la sección
\newbox\sectsavebox
\setbox\sectsavebox=\hbox{\boldmath\VT{xyz}}

%%%%%%%%%%%%%%%%%
\newbox\sectsavebox
\setbox\sectsavebox=\hbox{\boldmath\VT{xyz}}


%%%%%%%%%%%%%%%%%%%%%%%%%%%%%%%%%%%%%%%%%%%%%%%%%%%%%%%%%%%%%%%%%%%%%%%
%%%%%%%%%%%%%%%%%%%%%%%%%%%%%%%%%%%%%%%%%%%%%%%%%%%%%%%%%%%%%%%%%%%%%%%
%%%%%%%%%%%%%%%%%      INICIO DEL LIBRO     %%%%%%%%%%%%%%%%%%%%%%%%%%%
%%%%%%%%%%%%%%%%%%%%%%%%%%%%%%%%%%%%%%%%%%%%%%%%%%%%%%%%%%%%%%%%%%%%%%%
%%%%%%%%%%%%%%%%%%%%%%%%%%%%%%%%%%%%%%%%%%%%%%%%%%%%%%%%%%%%%%%%%%%%%%%

\begin{document}

\sloppy 
\decimalpoint
\nocite{*}
\setlength\tabcolsep{3pt}
\newpage\null\thispagestyle{empty}\newpage
\newpage\null\thispagestyle{empty}\newpage

%%%%%%%%%%%%% Titulo del libro: Primera carátula %%%%%%%%%%%%%%%%%%
\booktitle{Editorial UNSA: Plantilla libro}

%%%%%%%%%%%%%%%%%%%%%%%%%%% Autor/ Autores %%%%%%%%%%%%%%%%%%%%%%%%%
\authors{Autor 1 Autor 1 \\ Autor 2 Autor 2\\ Autor 3 Autor 3}


%%%%%%%%%%%%% Titulo del libro: Segunda carátula %%%%%%%%%%%%%%%%%%
\offprintinfo{Editorial UNSA: Plantilla libro}{Autor 1 Autor 1 \\ Autor 2 Autor 2\\ Autor 3 Autor 3}



\halftitlepage
\titlepage

%%%%%%%%%%%%% Página de créditos %%%%%%%%%%%%%%%%%%
\begin{copyrightpage}[Autor 1 \&  Autor 2 \&  Autor 3, Primera Edición]{}\\

Editado e Impreso en:
\copyright Editorial UNSA
Universidad Nacional de San Agustín de Arequipa
Calle Paucarpata, Puerta 5, Área de ingenierías
Teléfono: 215558
E-mail: editorial@unsa.edu.pe
Arequipa-Perú\\

\textit{Primera edici\'on, mes del 2022} 
Tiraje:   ejemplares\\



Hecho en el Depósito Legal en la Biblioteca Nacional del Perú N$^\circ$: 
ISBN: \\



Reservados todos los derecho. No se permite la reproducción total o 
parcial de esta obra, ni su incorporación a un sistema informático, 
ni su transmisión en cualquier forma o por cualquier medio (electrónico,
mecánico, fotocopia, grabación u otros) sin autorización previa y por 
escrito de los titulares del \textit{copyright}. La infracción de dichos derechos puede 
constituir un delito contra la propiedad intelectual.\\

IMPRESO EN PERÚ

\end{copyrightpage}

%%%%%%%%%%%%%%%% Dedicatoria %%%%%%%%%%%%%%%%%%
\begin{dedication}
LaTeX, we don’t need to do such copy and paste thing. LaTeX has different packages which automatically generates dummy text in our document. You can generate them with just a few lines of code.
\end{dedication}
%%%%%%%%%%%%% Contenido en breve %%%%%%%%%%%%%%%%%%
\contentsinbrief

%%%%%%%%%%%%% Índice General %%%%%%%%%%%%%%%%%%
\tableofcontents

%%%%%%%%%%%%% Índice de figuras %%%%%%%%%%%%%%%%%%
%\listoffigures

%%%%%%%%%%%%% Índice de tablas %%%%%%%%%%%%%%%%%%
%\listoftables

%%%%%%%%%%%%% Agradecimiento %%%%%%%%%%%%%%%%%%
\begin{acknowledgments}
A la Universidad Nacional San Agustín de Arequipa por la subvención brindada mediante el fondo concursable: Publicación de Libros y Textos Académicos, -----, contrato Nro. ----- UNSA. \\
\end{acknowledgments}

%%%%%%%%%%%%% Introducción %%%%%%%%%%%%%%%%%%
\begin{introduction}
\lipsum[2-4] %%Reemplazar por texto
\end{introduction}

%%%%%%%%%%%%% Presentación %%%%%%%%%%%%%%%%%%

\begin{presentation}
\lipsum[2-4] %%Reemplazar por texto
\end{presentation}

%%%%%%%%%%%%% Prefacio %%%%%%%%%%%%%%%%%%

\begin{preface}
\lipsum[2-4] %%Reemplazar por texto
\end{preface}


%%%%%%%%%%%%%%%%%%%%%%%%%%%%%%%%%%%%%%%%%%%%%%%%%%%
%%%%%%%%%%%%%%%%%%%%%%%%%%%%%%%%%%%%%%%%%%%%%%%%%%%
%Capitulo 1
%%%%%%%%%%%%%%%%%%%%%%%%%%%%%%%%%%%%%%%%%%%%%%%%%%%
%%%%%%%%%%%%%%%%%%%%%%%%%%%%%%%%%%%%%%%%%%%%%%%%%%%

\chapter{Lorem ipsum 1}

\lipsum[2-9]

\section{Lorem ipsum 1.1}

\lipsum[2-4]

%%%%%%%%%%%%%%%%%%%%%%% Insertar figuras
\begin{figure}
    \centering
    \includegraphics[width=6cm,height=6cm]{figura/unsa.png}
    \caption{Primera descripción.}
    \label{figura1}
\end{figure}


\lipsum[2-4]


\section{Lorem ipsum 1.2}


\subsection{Lorem ipsum 1.2.1}
\lipsum[2-4]
\subsection{Lorem ipsum 1.2.2}
\lipsum[2-4]
\subsection{Lorem ipsum 1.2.3}
\lipsum[2-4]

\subsubsection{Lorem ipsum 1.2.3.1:}
%%%% ejemplo de cita con el comando \cite{}
Lorem ipsum dolor sit amet, consectetur adipiscing elit \cite{Perez}. Etiam sollicitudin elit et urna commodo, a vulputate lectus bibendum. In hac habitasse platea dictumst. Integer ac elit ut quam imperdiet volutpat. Maecenas vitae condimentum orci. Vestibulum odio urna, porta ut pulvinar et, rutrum vel eros. Phasellus efficitur laoreet elit id mollis. Morbi convallis pretium dui, ac facilisis eros molestie ac. In at purus sed orci accumsan rutrum. Aenean lacinia nibh et ipsum iaculis, id facilisis tellus vulputate. Suspendisse pellentesque turpis nec risus bibendum ullamcorper. Morbi sed libero commodo, iaculis neque sit amet, consequat est. Vestibulum ante ipsum primis in faucibus orci luctus et ultrices posuere cubilia curae; Quisque ultrices lorem sit amet ex scelerisque, ac facilisis justo lacinia. Nulla sit amet justo nisl. Morbi nec blandit ante (Figura \ref{figura1}).

\subsubsection{Lorem ipsum 1.2.3.2:}
Lorem ipsum dolor sit amet, consectetur adipiscing elit. Etiam sollicitudin elit et urna commodo, a vulputate lectus bibendum. In hac habitasse platea dictumst. Integer ac elit ut quam imperdiet volutpat. Maecenas vitae condimentum orci. Vestibulum odio urna, porta ut pulvinar et, rutrum vel eros \cite{Nakamura}. Phasellus efficitur laoreet elit id mollis. Morbi convallis pretium dui, ac facilisis eros molestie ac (Tabla \ref{tabla1}). In at purus sed orci accumsan rutrum. Aenean lacinia nibh et ipsum iaculis, id facilisis tellus vulputate. Suspendisse pellentesque turpis nec risus bibendum ullamcorper. Morbi sed libero commodo, iaculis neque sit amet, consequat est. Vestibulum ante ipsum primis in faucibus orci luctus et ultrices posuere cubilia curae; Quisque ultrices lorem sit amet ex scelerisque, ac facilisis justo lacinia. Nulla sit amet justo nisl. Morbi nec blandit ante (Tabla \ref{tabla2}).

\subsubsection{Lorem ipsum 1.2.3.3:}
Lorem ipsum dolor sit amet, consectetur adipiscing elit. Etiam sollicitudin elit et urna commodo, a vulputate lectus bibendum. In hac habitasse platea dictumst. Integer ac elit ut quam imperdiet volutpat. Maecenas vitae condimentum orci. Vestibulum odio urna, porta ut pulvinar et, rutrum vel eros. Phasellus efficitur laoreet elit id mollis. Morbi convallis pretium dui, ac facilisis eros molestie ac. In at purus sed orci accumsan rutrum. Aenean lacinia nibh et ipsum iaculis, id facilisis tellus vulputate. Suspendisse pellentesque turpis nec risus bibendum ullamcorper. Morbi sed libero commodo, iaculis neque sit amet, consequat est. Vestibulum ante ipsum primis in faucibus orci luctus et ultrices posuere cubilia curae; Quisque ultrices lorem sit amet ex scelerisque, ac facilisis justo lacinia. Nulla sit amet justo nisl. Morbi nec blandit ante. (Figura \ref{figura2})


%%%%%%%%%%%%%%%%%%%%%%% Insertar figuras
\begin{figure}
    \centering
    \includegraphics[width=.5\textwidth]{figura/unsa.png}
    \caption{Segunda descripción.}
    \label{figura2}
\end{figure}

%%%%%%%%%%%%%%%%%%%%%%%%%%%%%%%%%%%%%%%%%%%%%%%%%%%
%%%%%%%%%%%%%%%%%%%%%%%%%%%%%%%%%%%%%%%%%%%%%%%%%%%
%Capitulo 2
%%%%%%%%%%%%%%%%%%%%%%%%%%%%%%%%%%%%%%%%%%%%%%%%%%%
%%%%%%%%%%%%%%%%%%%%%%%%%%%%%%%%%%%%%%%%%%%%%%%%%%%

\chapter{Lorem ipsum 2}

\lipsum[2-9]

%%%%%%%%%%%%%%%%%%%%%%% Insertar tablas
\begin{table}
\centering
\caption{Tabla descripción.}
\label{tabla1}
\begin{tabular}{llr}
\hline
\multicolumn{2}{c}{Item} &            \\ \cline{1-2}
Animal     & Description & Price (\$) \\ \hline
Gnat       & per gram    & 13.65      \\
           & each        & 0.01       \\
Gnu        & stuffed     & 92.50      \\
Emu        & stuffed     & 33.33      \\
Armadillo  & frozen      & 8.99       \\ \hline
\end{tabular}
\end{table}


%%%%%%%%%%%%%%%%%%%%%%% Insertar tablas
\begin{table}
\centering
\caption{Tabla 2 descripción.}
\label{tabla2}
\begin{tabular}{|l|l|r|}
\hline
\multicolumn{3}{|c|}{Item}          \\  \hline
Animal     & Description & Price (\$) \\ \hline
Gnat       & per gram    & 13.65      \\
           & each        & 0.01       \\
Gnu        & stuffed     & 92.50      \\
Emu        & stuffed     & 33.33      \\
Armadillo  & frozen      & 8.99       \\ \hline
\end{tabular}
\end{table}

%%%%%%%%%%%%%%%%%%%%%%% Insertar tablas
\begin{table}
\centering
\caption{Tabla 3 descripción.}
\label{tabla2}
\begin{tabular}{|cc|r}
\cline{1-2}
\multicolumn{2}{|c|}{\cellcolor[HTML]{6200C9}{\color[HTML]{FFFFFF} \textbf{Item}}}                                                 &                                                                                         \\ \hline
\rowcolor[HTML]{9698ED} 
\multicolumn{1}{|c|}{\cellcolor[HTML]{9698ED}{\color[HTML]{FFFFFF} \textbf{Animal}}} & {\color[HTML]{FFFFFF} \textbf{Description}} & \multicolumn{1}{c|}{\cellcolor[HTML]{9698ED}{\color[HTML]{FFFFFF} \textbf{Price (\$)}}} \\ \hline
\multicolumn{1}{|c|}{Gnat}                                                           & per gram                                    & \multicolumn{1}{r|}{13.65}                                                              \\ \hline
\multicolumn{1}{|c|}{}                                                               & each                                        & \multicolumn{1}{r|}{0.01}                                                               \\ \hline
\multicolumn{1}{|c|}{Gnu}                                                            & stuffed                                     & \multicolumn{1}{r|}{92.50}                                                              \\ \hline
\multicolumn{1}{|c|}{Emu}                                                            & stuffed                                     & \multicolumn{1}{r|}{33.33}                                                              \\ \hline
\multicolumn{1}{|c|}{Armadillo}                                                      & frozen                                      & \multicolumn{1}{r|}{8.99}                                                               \\ \hline
\end{tabular}
\end{table}


Lorem ipsum (Ecuación \ref{ecuacion1}) dolor sit amet, consectetur adipiscing elit. Etiam sollicitudin elit et urna commodo, a vulputate lectus bibendum. In hac habitasse platea dictumst. Integer ac elit ut quam imperdiet volutpat. Maecenas vitae condimentum orci. Vestibulum odio urna, porta ut pulvinar et, rutrum vel eros. Phasellus efficitur laoreet elit id mollis. Morbi convallis pretium dui, ac facilisis eros molestie ac. In at purus sed orci accumsan rutrum. Aenean lacinia nibh et ipsum iaculis, id facilisis tellus vulputate. Suspendisse pellentesque turpis nec risus bibendum ullamcorper. Morbi sed libero commodo, iaculis neque sit amet, consequat est. Vestibulum ante ipsum primis in faucibus orci luctus et ultrices posuere cubilia curae; Quisque ultrices lorem sit amet ex scelerisque, ac facilisis justo lacinia. Nulla sit amet justo nisl. Morbi nec blandit ante.

%%%%%%%%%%%%%%%%%%%%%%% Insertar ecuaciones
\begin{equation}\label{ecuacion1}
    x_{1/2} = - \dfrac{p}{2}{ \pm \sqrt {\left( \dfrac{p}{2} \right)^2 - q}}
\end{equation}

\begin{equation}\label{ecuacion2}
2\sqrt{2}\,,\quad 2^2\sqrt{2-\sqrt{2}}\,,\quad 2^3\sqrt{2-\sqrt{2+\sqrt{2}}}\,,
\quad 2^4\sqrt{2-\sqrt{2+\sqrt{2+\sqrt{2+\sqrt{2}}}}}\,,\;\ldots
\end{equation}

\begin{equation}\label{ecuacion3}
p_k(x)=\prod_{\substack{i=1\\i\ne k}}^n
\left(\frac{x-t_i}{t_k-t_i}\right)
\end{equation} 

\begin{equation}
    \overbrace{x+\underbrace{y+z}_{2} +w}^{4} 
\end{equation}



\section{Formulas químicas}
%%%%%%%%%%%%%%%%%%%%%%% Insertar formulas químicas

\begin{equation*}
    \chemfig{A*5(-B=C-D-E=)}
\end{equation*}

\begin{equation*}
\chemfig{A*6(-B=C(-CH_3)-D-E-F(=G)=)}
\end{equation*}

\begin{equation}
    \begin{array}{c}
        \chemfig{
            H_{3}\chemabove{N}{\scriptstyle+}- 
            CH(
                -[:90,,1,1]R^{1}
            )- 
            C(
                =[:-90]O
            )- 
            OH
            \qquad
            +
            \qquad
            H-
            N(
                -[:90]H
            )-
            CH(
                -[:90,,1,1]R^{2}
            )- 
            COO^{-}
        }
        \\
        \\
        \begin{tikzpicture}
            \node[anchor=south west,inner sep=0] at (0,0) {\includegraphics[scale=0.15]{figura/arrowsReacaoPeptideo.png}};
            \node at (-0.5,0.5) {H\textsubscript{2}O};
            \node at (1.8,0.5) {H\textsubscript{2}O};
        \end{tikzpicture}
        \\
        \\
        \chemfig{
            H_{3}\chemabove{N}{\scriptstyle+}- 
            CH(
                -[:90,,1,1]R^{1}
            )- 
            C(
                =[:-90]O
            )- 
            N(
                -[:90]H
            )-
            CH(
                -[:90,,1,1]R^{2}
            )- 
            COO^{-}
        }
    \end{array}
    \label{ecuacion5}
\end{equation}

\section{Mapas conceptuales}
%%%%%%%%%%%%%%%%%%%%%%% Insertar mapas conceptuales
\begin{tikzpicture}[mindmap, grow cyclic, every node/.style=concept, concept color=orange!40, 
	level 1/.append style={level distance=5cm,sibling angle=90},
	level 2/.append style={level distance=3cm,sibling angle=45},]
\node{ShareLaTeX Tutorial Videos}
child [concept color=blue!30]{ node {Beginners Series}
	child { node {First Document}}
	child { node {Sections and Paragraphs}}
	child { node {Mathematics}}
	child { node {Images}}
	child { node {bibliography}}
	child { node {Tables and Matrices}}
	child { node {Longer Documents}}
}
child [concept color=yellow!30]{ node {Thesis Series}
	child { node {Basic Structure}}
	child { node {Page Layout}}
	child { node {Figures, Subfigures and Tables}}
	child { node {Biblatex}}
	child { node {Title Page}}
}
child [concept color=teal!40] { node {Beamer Series}
	child { node {Getting Started}}
	child { node {Text, Pictures and Tables}}
	child { node {Blocks, Code and Hyperlinks}}
	child { node {Overlay Specifications}}
	child { node {Themes and Handouts}}
}
child [concept color=purple!50] { node {TikZ Series}
	child { node {Basic Drawing}}
	child { node {Geogebra}}
	child { node {Flow Charts}}
	child { node {Circuit Diagrams}}
	child [concept color=green!40]  { node {Mind Maps}}
};
\end{tikzpicture}

%%%%%%%%%%%%%%% Bibliografía %%%%%%%%%%%%%%%%%%

\startonoddpage
\markboth{BIBLIOGRAF\'IA}{BIBLIOGRAF\'IA}
\addcontentsline{toc}{chapter}{\fontsize{10}{1}\fontfamily{phv}\selectfont\textbf{BIBLIOGRAFÍA}}


%Muestra la bibliografía
\printbibliography

%%%%%%%%%%%%%%%%%%%%%%%%%%%%%%
%%%%%%%%%%%%%%%%%%%%%%%%%%%%%%

\end{document}






